\section{引言}

大语言模型的近期进展~\citep{gpt,gpt2,gpt3,gpt3.5,gpt4,llama,llama2}，为构建能够自主解决复杂任务的语言智能体带来了可能。开发 AI 智能体的常见做法，是将复杂任务分解为由提示和工具堆叠而成的 LLM 流水线~\citep{park2023generative,hong2023metagpt,zhou2023agents,chen2023agentverse,OpenAgents}。从某种意义上说，语言智能体可以被视为连接联结主义 AI（即智能体的 LLM 骨干）与符号主义 AI（即提示和工具组成的流水线）的 AI 系统；这也部分解释了它们在现实问题求解场景中的有效性。

然而，当前语言智能体的发展受制于这样一个事实：针对特定任务构建并定制语言智能体系统需要大量工程投入。具体而言，研究人员和开发者必须手工把复杂任务分解为对 LLM 更容易处理的子任务——本文称之为节点——然后再为特定节点精心设计提示与工具，包括 API 函数、知识库、记忆等。这一过程的复杂性，使当前语言智能体研究呈现出\textit{以模型为中心}或\textit{以工程为中心}的格局。这意味着，研究人员几乎无法像以\textit{数据为中心}训练神经网络那样，在数据集上手工调节或优化语言智能体。它既限制了手工编码语言智能体的鲁棒性和通用性，也使智能体适配新任务或新数据分布时需要付出大量工程成本。我们认为，从以工程为中心的语言智能体开发转向以数据为中心的学习，是语言智能体研究的重要一步。

为此，近期已有一些工作尝试自动优化语言智能体。例如，DSPy~\citep{khattab2023dspy}提出通过自举或在不同提示组件的组合空间中随机搜索，算法化地优化 LLM 提示；GPTSwarm~\citep{zhuge2024language}进一步提出以迭代优化过程处理 DSPy 所带来的组合优化挑战。Agent-pro~\citep{zhang2024agentpro}提出在竞争环境中优化与智能体内部策略相对应的提示组件。AgentOptimizer~\citep{zhang2024offline}则提出用精心设计的提示来优化函数。尽管这些方法在某些场景中有效，它们通常只优化智能体系统中的独立模块，例如某个特定节点的提示。因此，这些优化方法容易陷入孤立提示、工具或节点各自的局部最优，进而损害整个智能体系统的性能。这类似于早期训练神经网络时逐层独立优化的做法~\citep{hinton2006reducing}；如今看来，整体优化神经网络能够获得更好的性能几乎是不言自明的。我们认为，智能体优化同样如此：联合优化智能体内全部符号组件，是优化智能体的关键。

\begin{figure}[tp]
    \centering
    \includegraphics[width=0.72\linewidth]{figs/overview.pdf}
    \caption{\baby{}与神经网络联结主义学习之间的类比。}
    \label{fig:overview}
\end{figure}

本文提出一个用于训练语言智能体的\textit{\baby}框架。该框架受神经网络训练所采用的联结主义学习过程~\citep{hinton1990connectionist}启发。具体来说，我们在语言智能体与神经网络之间建立如下类比：智能体的流水线对应神经网络的计算图；流水线中的节点对应神经网络中的层；节点所用的提示和工具则对应一层的权重。由此，我们可以在智能体训练语境下，用基于语言的损失、梯度和权重来实现联结主义学习的主要组成部分，即反向传播和基于梯度的权重更新。具体而言，我们使用精心设计的提示流水线，实现智能体训练中的损失函数、反向传播和权重优化器。对于一个训练样本，框架首先执行“前向传播”（即运行智能体），并将每个节点的输入、输出、提示和工具使用情况存入一条“轨迹”。随后，我们用基于提示的损失函数评估结果，得到“语言损失”。之后，沿轨迹从最后一个节点向第一个节点反向传播语言损失，为每个节点中的符号组件生成文本分析与反思；我们称之为语言梯度。最后，再通过另一套精心设计的提示，依据语言梯度更新各节点中的全部符号组件，以及由节点及其连接构成的计算图。该方法也天然支持优化多智能体系统：既可以把不同智能体视为不同节点，也可以允许多个智能体在同一节点内采取行动。

\baby{}框架是一种模拟标准联结主义学习过程的智能体学习框架。不同于分别优化单个提示或工具的既有方法，\baby{}联合优化智能体系统中的全部符号组件，包括提示、工具，以及把它们堆叠成智能体系统的流水线。这种自顶向下的优化方案还使\baby{}能够“整体性”地优化智能体系统，避免各独立组件的局部最优。这使面向复杂现实问题的语言智能体能够有效地\textit{从数据中学习}，从而为语言智能体研究由以工程为中心转向以数据为中心打开了可能。此外，由于基于语言的损失函数在生成语言损失时不要求真实标签，我们的框架使语言智能体能够\textit{从经验中学习}，并在创建、部署到真实环境之后\textit{有意识地}更新自身的所有符号组件，从而实现“自演化智能体”\footnote{智能体也可以在真实环境中收集训练数据，并通过微调更新 LLM 骨干。这样，智能体中的全部组件都可以更新。我们将其留作未来工作。}。

作为概念验证，我们在标准 LLM 基准和复杂智能体任务上开展了一系列实验。结果表明，所提出的\baby{}框架能够有效优化、设计提示与工具，并通过从训练数据中学习来更新整体智能体流水线。我们开源了\baby{}框架的全部代码和提示，以推动未来对\textit{以数据为中心}的智能体学习研究。
